\section{PNEW-015: the actual hook wave and a tensor coefficient gate}
\label{sec:hook-covariant-gate}

This section concerns the actual trace series of K-023, not the formal
series of PNEW-014. It proves an exact residue identity and an
exponential wave estimate. Its tensor conclusion is an unconditional
dichotomy and a conditional coefficient bound; the required covariant
bound is not proved here.

Write $b_n(s)$ for the coefficient of the hook trace Hilbert series
$B_n(t)$ of K-023, and write
\[
       k_n(s)=b_n(s)+3c_n(s),\qquad c_n(s)\ge0.
                                                        \tag{PNEW-H1}
\]
Here $c_n(s)$ is the multiplicity of the two-dimensional irreducible
rational representation of $C_3$ in the actual rectangular multiplicity
space. The global quasipolynomials are indexed by their ordinary
integer argument $s$.

For clarity, global quasipolynomiality of $b_n$ requires no conjecture.
The positively graded reductive invariant ring
$\mathbb C[W\oplus\mathbb S_{21}W]^{SL_n}$ is normal and has rational
singularities. Thus it is Cohen--Macaulay and has negative $a$-invariant.
For a homogeneous system of parameters, its Hilbert numerator has
degree equal to the sum of the parameter degrees plus that negative
invariant. The Hilbert series is therefore proper. Partial fractions
at roots of unity give its quasipolynomial at every $s\ge0$.
The same established argument applies to $k_n$, and then to $c_n$
by (PNEW-H1). The standard rational-singularity input is Boutot's
reductive quotient theorem (1987, principal theorem, p.~65); the precise
graded implication used here is Kurano--Sato--Singh--Watanabe (2009),
Theorem~5.6, printed p.~3262.
No exact formula for either ring's $a$-invariant is needed.

\subsection*{The exact residue identity, with both grading characters}

Put $M=2^{n-1}$ and fix a primitive $M$th root $\eta$. Use the
diagonal lift with eigenvalues
\[
 z_i=\eta^{(-2)^{i-1}}\ (1\le i<n),\qquad z_n=1.
\]
They are distinct and satisfy $z_i^2=z_{i+1}^{-1}$ for $i<n$.
Define the nonzero products
\[
 \begin{split}
 L&=\prod_{i<n}(1-z_i),\qquad
 T=\prod_{i<n}(1-z_i^3),\qquad
 R=\prod_{j=2}^n(1-z_j/z_1),\\
 D&=\prod_{i<j<k}(1-z_i z_j z_k),\qquad
 N=\prod_{i\ne j}(1-z_i/z_j),\\
 E&=\prod_{\substack{i\ne j\\z_i^2z_j\ne1}}(1-z_i^2z_j).
 \end{split}                                             \tag{PNEW-H2}
\]
The complete fixed-weight classification previously proved in P-F30
and K-024 ensures that the factors in $D$ are nonzero and that the
excluded factors in $E$ are exactly $j=i+1$, $i<n$.

Let $C_{\rm cub}(\eta)$ be the cubic chain's density-one local
coefficient, and $C_{\rm hook}(\eta)$ the hook chain coefficient.
These are evaluated at their respective grading roots. Specifically,
if $g_n=\gcd(n,3)$ and
$S=\sum_{i=0}^{n-2}(-2)^i$, the cubic grading root is
$\eta^{3S/g_n}$ and the hook grading root is $\eta^S$.
Since $3S\equiv1\pmod M$, both are primitive.

\begin{proposition}[PNEW-015a]
The exact moving-factor coefficients satisfy
\[
 C_{\rm cub}(\eta)=\frac{R}{LD},\qquad
 C_{\rm hook}(\eta)=\frac{RT}{L^2D^2}
                  =C_{\rm cub}(\eta)^2\frac{T}{R}.
                                                        \tag{PNEW-H3}
\]
\end{proposition}
\begin{proof}
The roots of the special-linear Lie algebra give $N$. The moving
cubic weights consist of triple weights, repeated weights and
distinct-index weights, with products $T,E,D$. The moving hook
representation $W\oplus\mathbb S_{21}W$ has products $L,E,D^2$:
each distinct-index weight in $\mathbb S_{21}$ has multiplicity two.
Thus the two coefficients are $N/(TED)$ and $N/(LED^2)$.

In the repeated-weight product, replace $z_i^2z_j$ by
$z_j/z_{i+1}$ for $i<n$. These are all root factors whose denominator
index is at least two, except the predecessor factors
$1-z_i/z_{i+1}=1-z_i^3$. The denominator-index-one root factors
form $R$. The terms $i=n$ contribute $L$. Consequently
$E=NL/(TR)$. Substitution proves every equality in (PNEW-H3).
The density-one conventions are exactly those of P-F28 and K-024;
no additional finite-stabilizer factor is inserted.
\end{proof}

\subsection*{Dominance of the actual auxiliary wave}

Let $P_n(a)$ and $Q_n$ be the cubic amplitudes and relative remainder
from the already reviewed quantitative refinement. Thus
$|C_{\rm cub}(\exp(2\pi ia/M))|=P_n(a)$, and the selected lift
indices are $a=\pm3^{-1}\pmod M$. Write $Q_n^{\rm hook}$ for the
analogous hook remainder, and let
\[
 A_n=\frac{|C_{\rm hook}(\exp(2\pi i3^{-1}/M))|}{M}.
\]

\begin{proposition}[PNEW-015b]
For every $n\ge79004$,
\[
 Q_n^{\rm hook}<\tfrac14,\qquad
 A_n\ge2^{n^4/12-5n^3}.                                \tag{PNEW-H4}
\]
If $W_n^b(x)$ is the full primitive-$M$ wave of $b_n$, every nine
consecutive integer arguments contain one $x$ with
\[
                       2W_n^b(x)\ge A_n/2.
                                                        \tag{PNEW-H5}
\]
In particular $p(b_n)>2^{n/8}$ and $s(b_n)\ge2^{n/8}$ under the
same strict constituent definitions.
\end{proposition}
\begin{proof}
Every nonzero dyadic factor has modulus between $4/M$ and two.
With $C=(n-1)(n-2)$ this gives $2^{-C}\le|T/R|\le2^C$.
Proposition PNEW-015a therefore implies
\[
                  Q_n^{\rm hook}\le2^{2C}Q_n^2.
                                                        \tag{PNEW-H6}
\]
Use the already proved recurrence
$Q_j\le2^{1-2j^2/7}Q_{j-1}+2\cdot2^{-j^3/12}$ and
$Q_{n-4}<1/4$. Four iterations, bounding all contraction factors
by one in the error terms, give
\[
 Q_n\le2^{2-8(n-3)^2/7}+2^{3-(n-3)^3/12}.
\]
Using $(a+b)^2\le2a^2+2b^2$ in (PNEW-H6), the remainder is at most
\[
 2^{2C+5-16(n-3)^2/7}
          +2^{2C+7-(n-3)^3/6}<\tfrac14.                 \tag{PNEW-H7}
\]
Both exponents are below $-3$ at $n=79004$ and decrease for $n\ge14$;
the retained exact rational certificate checks the endpoint margins.
The cubic amplitude bound gives
$P_n(3^{-1})\ge2^{n^4/24-2n^3}$. Squaring, multiplying by $2^{-C}$
and dividing by $M$ proves the second part of (PNEW-H4), since
$C+n-1=(n-1)^2\le n^3$.

The selected hook grading index is $3^{-1}S=9^{-1}\pmod M$.
Its nine successive phases differ from nine equally spaced phases
by at most $2\pi/M$. One of their cosines is at least
$\cos(\pi/9)-2\pi/M>1/2$ for $n\ge6$.
The selected conjugate pair therefore has a value at least $A_n$;
subtracting the remaining absolute mass, less than $A_n/4$, proves
the weaker convenient bound (PNEW-H5).

The auxiliary representation has dimension $N=(n^3+2n)/3$.
Counting all monomials of ordinary degree at most $nx$ bounds its
weighted-degree-$nx$ invariant dimension by $\binom{nx+N}{N}$.
Take the nine-argument window beginning at $\lceil2^{n/8}\rceil$.
Here $x\le2^{n/8+1}$ and $nx\ge N$. For $n\ge79004$,
$2+\log_2n\le n/32$, so
\[
 \log_2 b_n(x)\le N(n/8+2+\log_2n)
       \le5n^4/96+5n^2/48
       < n^4/12-5n^3-1.
\]
The logarithmic bound is interpreted as an upper bound on the
ambient binomial, so also covers $b_n(x)=0$. Thus $2W_n^b(x)>b_n(x)$.
K-024 excludes every other pole whose two-part is at least $M$.
The constituent half-period identity of P-F56 consequently applies
to this actual auxiliary series and gives a negative nonzero
constituent at $x$. This proves the two strict-index conclusions.
\end{proof}

\subsection*{A comparison that allows higher tensor two-parts}

For $n\ge4$ put $D_n=n^3-3n^2+3$. Every constituent of $c_n$ has
degree at most $D_n-1$: along its own residue, $0\le c_n(s)\le k_n(s)/3$,
and the latter has growth degree $D_n-1$. If
$c_{n,j}(X)=\sum_{d=0}^{D_n-1}a_{n,j,d}X^d$, define
\[
 U_c^+(n)=\max_j\sum_d\max(a_{n,j,d},0).
                                                        \tag{PNEW-H8}
\]
Enlarging the quasipolynomial period only repeats these polynomials,
so this quantity is independent of the chosen common period.

\begin{theorem}[PNEW-015c: an actual tensor coefficient dichotomy]
For $n\ge79004$, either
\[
       p(k_n)>2^{n/32}\quad\hbox{and}\quad
       s(k_n)\ge2^{n/32},
                                                        \tag{PNEW-H9}
\]
or
\[
       U_c^+(n)\ge2^{(5/96)n^4-6n^3-4}.
                                                        \tag{PNEW-H10}
\]
In particular, if a constant $C\ge1$ satisfies
\[
 U_c^+(n)\le2^{C n^3\log_2(n+2)},\qquad
 C\log_2(n+2)+8\le5n/96,                              \tag{PNEW-H11}
\]
then both strict tensor bounds (PNEW-H9) hold. The coefficient
estimate in (PNEW-H11) is an unproved gate for the actual covariant
multiplicity, not a consequence of (PNEW-H1).
\end{theorem}
\begin{proof}
Choose any common period $P=2^vQ$ of $k_n,b_n,c_n$, with $Q$ odd.
Since $b_n$ has a primitive-$M$ pole, $v\ge n-1$. Put
$h=(M/2)Q$. This need not be half of $P$.
All lower-two-part poles of $B_n$ have orders dividing $h$,
while its primitive-$M$ terms change sign under a shift by $h$.
K-024 excludes all other terms with two-part at least $M$.
Consequently, for every $j$ and every argument $X$,
\[
 b_{n,j+h}(X)=b_{n,j}(X)-2W_{n,j}^b(X),                \tag{PNEW-H12}
\]
where the primitive wave is the constant
$W_{n,j}^b=M^{-1}\sum_{\xi\ {
m primitive}\ M}
C_{\rm hook}(\xi)\xi^{-j}$, using the grading-root convention.
At an integer $x$ with $j\equiv x\pmod P$, this wave is $W_n^b(x)$.
No condition on larger two-parts in the tensor or covariant series
was used.

Choose $x$ in the nine-argument window beginning at
$\lceil2^{n/32}\rceil$ for which (PNEW-H5) holds, and put $j=x\bmod P$.
The actual dimension bound gives
\[
 b_n(x)\le k_n(x)\le\binom{nx+n^3-1}{n^3-1}
          \le2^{n^3(n/32+2+\log_2n)}<A_n/4.           \tag{PNEW-H13}
\]
For the last inequality use $2+\log_2n\le n/32$ and
$n^4/16<n^4/12-5n^3-2$, valid in the stated range.
If either bound in (PNEW-H9) fails, the corresponding definition
forces every tensor constituent to be nonnegative at this $x$.
Using (PNEW-H1) and (PNEW-H12), therefore,
\[
 0\le k_{n,j+h}(x)
     =b_n(x)-2W_n^b(x)+3c_{n,j+h}(x).
\]
It follows that $c_{n,j+h}(x)>A_n/12$.
For $x\ge1$, discarding negative coefficients bounds this value by
$U_c^+(n)x^{D_n-1}$. Since $x\le2^{n/32+1}$, we obtain
\[
 \begin{split}
 \log_2 U_c^+(n)
 &>n^4/12-5n^3-\log_2 12-(D_n-1)(n/32+1)\\
 &\ge(5/96)n^4-6n^3-4.
 \end{split}
\]
This proves the dichotomy. The second inequality in (PNEW-H11)
makes its exponent bound strictly smaller than the required lower
bound, since $n^3>2$. Hence (PNEW-H10) is impossible in that case,
and both strict-index conclusions follow.
\end{proof}

\paragraph{What remains.}
The result establishes actual auxiliary exponential indices and a
specific positive-coefficient obstruction in the actual cyclic
covariant series if tensor indices are small. It does not bound
$U_c^+$, determine the exact tensor dyadic pole, or prove tensor
PH2/SH(a) false unconditionally. Attained nonnegativity of $c_n(s)$
does not imply positivity of its constituent polynomials away from
their residues. No such inference occurs in the proof.

\paragraph{Primary references and reading scope.}
Boutot, J.-F. (1987). Singularit\'es rationnelles et quotients par les
groupes r\'eductifs. \emph{Inventiones Mathematicae, 88}(1), 65--68.
\url{https://doi.org/10.1007/BF01405091}.
The input is the principal theorem on p.~65, applied to a polynomial
ring with its reductive $SL_n$ action. The publisher's primary
bibliographic record was checked in this branch; the original proof
was not newly read here, and remains a standard cited theorem input.

Kurano, K., Sato, E.-i., Singh, A.~K., \& Watanabe, K.-i. (2009).
Multigraded rings, diagonal subalgebras, and rational singularities.
\emph{Journal of Algebra, 322}(9), 3248--3267.
\url{https://doi.org/10.1016/j.jalgebra.2008.11.014}.
Theorem~5.6 and its complete proof on printed p.~3262, together with
the preceding Lemma~5.4 used in that proof, were read in the author's
primary PDF:
\url{https://math.utah.edu/~singh/publications/kurano_sato_singh_watanabe.pdf}.
Only the implication from rational singularities to Cohen--Macaulayness
and negative $a$-invariant is used above.
